\section{Introduction}\label{sec:introduction}
A finite instruction can prescribe a volume-preserving affine map on a
box. A fluid realization must do more: it must move every point of that
box, preserve incompressibility outside it, and leave the other instructions
available. Source boxes and destination boxes can overlap. Moreover, a
particle used to report the answer must avoid the observation region
throughout every nonhalting transition. We study these requirements together.

Our geometric result is Theorem~\ref{lem:bb-balanced-box-routing}.
It extends any finite collection of positive diagonal determinant-one maps
between separately disjoint rational solid boxes to the time-one map of
an effective smooth compactly supported incompressible velocity. The maps
hold on neighborhoods of the entire closed boxes. The proof separates
centers, changes shapes, evacuates the sources to storage, and delivers
boxes along paths avoiding everything still occupied. Positive margins make
the construction effective and permit additional observation guards.

The first computational application, Theorem~\ref{thm:bb-balanced-three-stack},
uses all three coordinates for symbolic data. Three stacks retain a tape
and its execution history. Every local instruction preserves the total
length of the prefixes it replaces, so its affine map has determinant one.
A fixed particle, initially in a fluid at rest, enters the fixed cube
$(-1,2)^3$ exactly when the encoded machine halts. The force is smooth,
compactly supported in space and periodic after a one-unit loading interval.
Its prescription uses the finite rule table, not the executed computation.
This application completes the main construction before we vary
initialization, history guards and observers.

The passage from a symbolic instruction to a fluid observation begins
with retained history, which makes the instruction injective on its
entire domain of tapes, not just along the chosen computation. A rule
specifies finite prefixes and leaves their infinite tails free; gapped
radix expansions place the codes in separately disjoint closed source
and target boxes. Prefix replacement extends to an affine map of each
whole box, which smooth incompressible transport then realizes.
An additional geometric bound must exclude unintended visits to
the observer between successive rule times. Finally, the prescribed
velocity determines a body force, and an energy comparison identifies
the resulting solution in the declared class. Each later variant changes
one or more of these steps while retaining this separation of tasks.

The remaining constructions answer more specific questions. Two tape
coordinates can be thickened into solid boxes, or their planar scale changes
can be compensated in the normal direction. A temporary mark, a persistent
mark, an explicit frontier and a boundary between occupied and unused
history give different partial instruction maps. We prove their inverses
on their full domains because this is what separates the target boxes.
A bounded observation box permits horizontal motion above it; a slab needs
a coordinate bound throughout that motion. Storage and obstacle detours
supply further alternatives. Loading once gives eventual periodicity,
whereas placing the initial code at the fixed label or repeating a loading
branch gives periodicity from time zero. Onto slow clocks give separate
decaying forces with derivative norms square-integrable in time.

\subsection{Prior work and the role of forcing}
The computational starting point is Turing's finite machine model
\cite{Turing1936}. Landauer's discussion of retained intermediate information
\cite[Section~3]{Landauer1961} and Bennett's explicit history recording
\cite[Eqs.~(9)--(11)]{Bennett1973} explain how an irreversible computation
can be embedded in an invertible mechanism. We use that principle, but
prove the finite guarded tables and their complete inverses here. No
complexity bound or history-erasure procedure from those works is asserted
for these tables.

Moore's generalized shifts turn symbolic updates into affine operations
on positional expansions \cite{Moore1990,Moore1991}. Our gapped expansions
serve the same purpose. Moore also realizes invertible generalized shifts
by smooth planar diffeomorphisms and three-dimensional flows
\cite[Section~6, Theorem~12 and its corollary]{Moore1991}.
Cardona, Miranda, Peralta-Salas and Presas extend the affine maps on
neighborhoods of separated planar blocks to an area-preserving disk
diffeomorphism, using contraction, transport, expansion and Moser's area
correction \cite[Proposition~5.1]{CardonaMirandaPeraltaSalasPresas2021}.
This is a geometric predecessor of the routing problem here. Our local
curl formulas implement positive diagonal determinant-one motions in
three dimensions directly; the proof supplies effective support and
neighborhood bounds and controls observation at intermediate times.

Computational universality has been established for stationary Euler
flows with adapted geometry \cite{CardonaMirandaPeraltaSalasPresas2021}.
Cardona, Miranda and Peralta-Salas also construct universal Euclidean
Beltrami flows, with infinite energy, and robust simulations of
tape-bounded machines by unforced Navier--Stokes flows on a flat torus
\cite[arXiv version~3, Theorem~1 and Remark~29]{CardonaMirandaPeraltaSalas2023Beltrami}.
Dyhr, Gonz\'alez-Prieto, Miranda and Peralta-Salas obtain stationary
unforced Navier--Stokes universality using an adapted metric and the
Hodge Laplacian
\cite[Theorem~A]{DyhrGonzalezPrietoMirandaPeraltaSalas2026}.

Here the ambient metric is fixed and flat, the viscosity is
prescribed, and the initial velocity is zero. The machine is placed in an
external body force. Given an explicit solenoidal velocity $U$, the identity
$f=U_t+(U\cdot\nabla)U-\nu\Delta U$ supplies that force. The mathematical
work is the effective construction of $U$ and its all-time observation,
followed by uniqueness in the stated comparison class. We do not assert a
regularity theorem for arbitrary fluid data.

The companion \cite{OpenAI379Shear} gives an entry construction using
solenoidal shears, and \cite{OpenAI379Sheets} treats prefix instructions
and other planar or normal-compensated realizations. They explain adjacent
geometric choices; the proofs of the present full-box mechanisms are given
here. The final input-offset application explicitly imports the planar
processor and spatial clock of \cite{OpenAI379StationaryA2}, with the exact
parameters and conclusions stated in Section~\ref{bc:input-offset-section}.
The companion \cite{OpenAI379EulerianC} detects computation from an
integral of the velocity field on $\R^2\times\T$ via an
advection--diffusion estimate. The end of
Section~\ref{sec:bb-initialization} compares that event with the path of
one material particle studied here.

\subsection{Equation, effectiveness and reading path}
Throughout $\nu>0$ is fixed, $\T=\R/\Z$, and the fluid equation on
$\R^3$ or the flat unit torus is
\begin{equation}\label{eq:box-NS}
 u_t+(u\cdot\nabla)u=-\nabla p+\nu\Delta u+f,
 \qquad \divg u=0,\qquad u(0)=0.
\end{equation}
A material label $a$ has path $\dot\Phi_t(a)=u(t,\Phi_t(a))$,
$\Phi_0(a)=a$. We also write $X(t;a)=\Phi_t(a)$. The smooth constructed
velocities have global material flows. On the torus pressure is normalized
to have mean zero. Whole-space pressure assumptions are specified in
Section~\ref{sec:model} and in the individual statements. Unless explicitly
projected, the force is a general body force and the constructed pressure
is zero.

An effective field is a finite program that evaluates every requested
mixed derivative at computably supplied arguments to any prescribed
rational error and computes the bounds used in the proof. The finite
machine and its finite input determine the coefficients, cutoffs and
routing schedule. Evaluating the field cannot first compute the selected
machine run. All algorithmic statements concern computable $\nu$;
the same formulas are valid for arbitrary positive $\nu$, with evaluation
relative to that parameter. The encoded trajectories are exact real
trajectories. No uniform precision tolerance for an unbounded computation
is asserted.

Section~\ref{sec:model} states the analytic comparison tools.
Section~\ref{sec:bb-boxes} proves the box motions, and
Section~\ref{sec:bb-balanced-three-stack} gives the balanced-stack application.
Sections~\ref{sec:bb-frontiers} and~\ref{ba:codes} establish reusable
recorders and full closed-rectangle coding; Section~\ref{ba:geometry}
supplies the common interpolation formulas. The subsequent constructions
can be read according to the requirement being changed:
\begin{center}
\begin{tabular}{@{}p{.24\textwidth}p{.71\textwidth}@{}}
\toprule
Choice & Where it is developed \\
\midrule
Initialization and time periodicity &
Section~\ref{sec:bb-initialization} compares a separate loader, direct
placement of the initial code, and a loading branch repeated in every
period. An untargeted start square permits the last choice. \\[4pt]
History domains &
Sections~\ref{ba:marked-routes} and~\ref{sec:bb-guarded} distinguish
persistent marks, fresh history and neighboring-cell guards. Their
inverses on arbitrary allowed tapes determine the separated image boxes.
\\[4pt]
Observer geometry &
Section~\ref{ba:ordered-routes} treats bounded boxes, slabs and
half-spaces. Sections~\ref{sec:solid-routing} and~\ref{sec:torus-routing}
develop storage, detours and fixed torus charts with the corresponding
intermediate-time bounds. \\[4pt]
Decay and stationarity &
The logarithmic clock in Section~\ref{ba:euclidean} produces separate
decaying forces. Section~\ref{sec:clocks} treats further profile choices,
including the imported input-offset processor. \\
\bottomrule
\end{tabular}
\end{center}
The five whole-space comparison classes are stated in
Lemma~\ref{lem:b-comparison}. Each alternative fixes its own label,
observer, support, pressure class and time dependence; these are separate
constructions rather than simultaneous guarantees for one force.
